% ECONOMICS_PROBLEM: {"id": "D82-64", "title": "Manipulation-resistant learning", "jel_code": "D82", "source_id": "OP-0287"}
% !TeX program = xelatex
\documentclass[11pt,a4paper]{article}
\usepackage[margin=23mm]{geometry}
\usepackage{fontspec}
\setmainfont{Latin Modern Roman}
\usepackage{amsmath,amssymb,mathtools}
\usepackage{xcolor,enumitem,booktabs,longtable,array,xurl,hyperref}
\definecolor{muted}{HTML}{53636C}
\hypersetup{colorlinks=true,urlcolor=blue,linkcolor=blue}
\setlength{\parindent}{0pt}
\setlength{\parskip}{5pt}
\newcommand{\R}{\mathbb R}
\newcommand{\E}{\mathbb E}
\newcommand{\N}{\mathbb N}
\newcommand{\ind}{\mathbf 1}
\newcommand{\Var}{\operatorname{Var}}
\newcommand{\Cov}{\operatorname{Cov}}
\newcommand{\supp}{\operatorname{supp}}
\DeclareMathOperator*{\argmin}{arg\,min}
\DeclareMathOperator*{\argmax}{arg\,max}
\newcommand{\entrymeta}[3]{{\sffamily\small\color{muted}\textbf{\texttt{#1}}\quad|\quad #2\par #3\par}}
\newcommand{\Problemref}[1]{\texttt{#1}}
\newcommand{\JELref}[2]{\texttt{#2} (JEL \texttt{#1})}
\begin{document}
\section*{Manipulation-resistant learning}
\textbf{Problem ID:} \texttt{D82-64}\quad \textbf{JEL:} \texttt{D82}\par
% BEGIN ECONOMICS STATEMENT
\label{problem:OP-0287}
{\sffamily\small\color{muted}Original subject group: Game theory and strategic interaction\par}
\label{definitions:problem:30007569}
\textbf{Audit classification:} Published research agenda. Okumura’s February 2026 v3 conclusion explicitly proposes multiple learning agents. The two-buyer benchmark and loss frontier are editorial, with no claim that their exact unsolved status was established.
\subsection*{Definitions and precise research target}
\textbf{Model and research target.} Consider a repeated market with one seller and two buyers. Buyer \(i\) privately observes a persistent type \(\theta_i\) in a known finite set \(\Theta_i\subseteq[0,1]\); the seller knows a full-support joint prior \(F\) from a specified class \(\mathcal P\), but not realizations. Each period buyers choose finite actions, and a seller policy maps their actions into a feasible allocation of one good and nonnegative payments. Buyers observe only their own realized allocations and payments and may opt out for zero payoff. Before interaction each buyer commits to a learning algorithm \(L_i\). Define stationary learning loss \(R_T(L)\) against the best fixed action when seller policies and other buyers’ actions are independently drawn each period from fixed unknown distributions. Against adaptive sellers, let \(\Sigma^{\eta_T}(L,F)\) be seller responses within \(\eta_T\) of maximal expected average revenue to the committed algorithm profile. For a static revenue-optimal Bayesian mechanism \(M_F^*\), define extraction loss
\[
 S_T(L)=\sup_{F\in\mathcal P}\sup_{\sigma\in\Sigma^{\eta_T}(L,F)}\max_{i,\theta_i}
 \left(U_i(M_F^*\mid\theta_i)-E_{L,\sigma}[T^{-1}\textstyle\sum_t u_{it}\mid\theta_i]\right).
\]
Characterize the attainable tradeoff between sublinear \(R_T\) and vanishing or bounded \(S_T\), including necessary information and participation assumptions. Define buyer utility \(u_{it}=\theta_i x_{it}-p_{it}\) and ensure the static mechanism uses identical feasibility constraints.

\emph{Definitions and maintained assumptions (editorial unless a source definition is named).} Types \(\theta_i\) are drawn once and persist; a buyer’s value for allocation probability \(x_{it}\) is \(\theta_i x_{it}\), and utility is \(u_{it}=\theta_i x_{it}-p_{it}\). Fix finite actions including opt-out, \(x_{1t}+x_{2t}\le1\), and a payment bound \(0\le p_{it}\le\bar p\), with opt-out enforcing \(x_{it}=p_{it}=0\). Specify what the seller and buyers observe and whether allocations and payments are publicly or privately observed. Stationary regret compares realized learning payoff with the expected payoff of the best fixed action under the stipulated stationary environment, conditional on each type. Revenue-best responses need not attain a maximum; use \(\Sigma^{\eta_T}(L,F)\), policies within \(\eta_T>0\) of the supremum expected average revenue, with \(\eta_T\downarrow0\). Replace \(\Sigma^*\) by this nonempty approximate set in \(S_T\). For \(M_F^*\), fix a canonical revenue-optimal static Bayesian mechanism and equilibrium selection with identical participation and payment bounds; different optimal mechanisms may leave different type utilities. Take positive-part extraction losses when a nonnegative regret frontier is desired.
\subsection*{Variants and assumption changes}
\begin{enumerate}
\item \textbf{Independent types (editorial).} Restrict \(F\) to full-support independent finite types and bandit payoff feedback. Seek algorithms combining sublinear stationary loss with bounded extraction loss uniformly over that class.
\item \textbf{Correlated and public feedback (editorial).} Admit correlation and public action/allocation observations. Determine whether another buyer’s learning reveals type information that invalidates the single-agent guarantee, and formulate the appropriate interim-utility benchmark.
\end{enumerate}
\subsection*{References and source audit}
\begin{enumerate}
\item §8 Conclusion, printed p.20 (v3 February2026). Supports the stated research area; exact displayed specification is editorial. \url{https://arxiv.org/pdf/2505.05341}
\end{enumerate}
\begin{enumerate}\raggedright
\item\hypertarget{definitions:source:30007569:1}{} Kyohei Okumura. 2026. \emph{Robust Learning with Private Information}. §8 Conclusion, printed p.20, final paragraph (version3, 10February2026).
\url{https://arxiv.org/pdf/2505.05341}.
\end{enumerate}

\subsection*{Common conventions: Source mathematical conventions}

These conventions apply to each entry unless explicitly replaced.

\textbf{Models and identification.} A primitive model \(m\) specifies all probability laws, feasible choices, information, equilibrium and measurement rules used by its target. The admissible class \(\mathcal M\) is fixed independently of outcomes. If \(P_O(m)\) is its observable probability law and \(\tau(m)\) the target, its identified set at observable law \(P\) is
\[
 \mathcal I_\tau(P)=\{\tau(m):m\in\mathcal M,\ P_O(m)=P\}.
\]
Point identification means that this set is a singleton. An empty set signals incompatibility of data and assumptions. Coordinate bounds need not describe the joint set. Bounds are sharp only if every retained target is attainable or arbitrarily approachable by compatible models. Finite bounds require explicit restrictions on unobserved quantities; unrestricted confounding, hidden wealth, extrapolation or arbitrary equilibrium selection can yield uninformative sets.

\textbf{Causal interventions.} A potential outcome \(Y_i(p)\) is the outcome under a complete policy regime \(p\), including timing, financing, information and market spillovers. The target population and units are fixed before treatment. With interference, use the complete assignment vector or a justified exposure mapping; individual no-interference notation alone cannot identify an economy-wide intervention. A difference of mean potential outcomes does not require the cross-world joint distribution. Conditioning on post-treatment migration, survival, enrollment or employment changes the estimand and needs separate justification.

\textbf{Statistical statements.} An estimator, confidence set, forecast or test specifies a sampling law, model class and loss/coverage criterion. Uniform \((1-\alpha)\)-coverage means \(\inf_{m\in\mathcal M}\Pr_m\{\tau(m)\in C_n\}\geq1-\alpha\), or an explicitly asymptotic version. The whole parameter space trivially has coverage, so informativeness requires a stated width, risk or power objective. A forecasting improvement is relative to a named benchmark, information set and evaluation population.

\textbf{Equilibrium and optimization.} An equilibrium comprises feasible plans, optimality, beliefs where needed, budget constraints and all market/resource clearing. The primitives or a stated benchmark define each correspondence \(\mathcal E(p;m)\). Nonemptiness is assumed or proved, never inferred just from compact policy sets. A selected-equilibrium target uses a declared selection rule; a robust target quantifies over all permitted equilibria. Use suprema or infima unless attainment is established. An \(\varepsilon\)-optimal policy is feasible and within \(\varepsilon\) of the finite optimum value. Report an empty feasible set separately.

\textbf{Welfare and accounting.} Normative utility, interpersonal normalization, social weights, numeraire, discounting and the use of public receipts are primitives. Behavioral utility need not coincide with the welfare criterion. Household welfare includes ownership income and rents once; adding profits again to compensating variation generally double-counts them. A monetary equivalent is defined by an explicit transfer and price convention, with monotonicity and range conditions. Average effects, mechanism contributions, transfers and welfare losses are different targets. Infinite sums require convergence and dynamic feasible plans require terminal or no-Ponzi conditions.

\textbf{Source versions.} A working-paper date, revision date and journal publication year are distinguished. A DOI for a journal version is not automatically the DOI of an earlier manuscript. Locators refer to the stated version and printed pages where specified. A metadata-only check does not verify a claim about a full-text conclusion. Every entry's source subsection narrows the provenance claim accordingly.
% END ECONOMICS STATEMENT
\end{document}
